% !PRINTED 305
\weberpage{305}

Wenn wir zwei auf einander folgende Glieder dieser Doppelreihe, wie
\[
\begin{array}{l}
f_\nu,\ f_{\nu+1} \\
F_\nu,\ F_{\nu+1}
\end{array}
\]
für einen Werth von $x$ betrachten, für den keine der vier Functionen
verschwindet, so sind in Bezug auf die Zeichenfolge die nachstehenden
Fälle möglich:
\[
\text{a)}\quad
\begin{array}{c}+\ +\\+\ +\end{array},\
\begin{array}{c}+\ +\\-\ -\end{array},\
\begin{array}{c}-\ -\\+\ +\end{array},\
\begin{array}{c}-\ -\\-\ -\end{array},\quad PP,
\]
\[
\text{b)}\quad
\begin{array}{c}+\ -\\+\ -\end{array},\
\begin{array}{c}+\ -\\-\ +\end{array},\
\begin{array}{c}-\ +\\+\ -\end{array},\
\begin{array}{c}-\ +\\-\ +\end{array},\quad VV,
\]
\[
\text{c)}\quad
\begin{array}{c}+\ +\\+\ -\end{array},\
\begin{array}{c}+\ +\\-\ +\end{array},\
\begin{array}{c}-\ -\\+\ -\end{array},\
\begin{array}{c}-\ -\\-\ +\end{array},\quad PV,
\]
\[
\text{d)}\quad
\begin{array}{c}+\ -\\+\ +\end{array},\
\begin{array}{c}+\ -\\-\ -\end{array},\
\begin{array}{c}-\ +\\+\ +\end{array},\
\begin{array}{c}-\ +\\-\ -\end{array},\quad VP.
\]
Dies Verhalten bezeichnen wir mit folgenden Ausdrücken und Symbolen:

a) Folge-Folge (Permanenz-Permanenz) $PP$,

b) Wechsel-Wechsel (Variation-Variation) $VV$,

c) Folge-Wechsel (Permanenz-Variation) $PV$,

d) Wechsel-Folge (Variation-Permanenz) $VP$.

Die zwei Sätze, die wir beweisen wollen, lauten:

\gesperrt{II. Die Anzahl der zwischen $\alpha$ und $\beta$ gelegenen
Wurzeln von $f(x)$ ist entweder genau so gross oder um eine gerade Zahl
kleiner, als die Zahl der beim Uebergang von $\alpha$ zu $\beta$ in der
Doppelreihe (7) verlorenen Wechsel-Folgen.}

\gesperrt{III. Die Anzahl der zwischen $\alpha$ und $\beta$ gelegenen
Wurzeln von $f(x)$ ist entweder genau so gross oder um eine gerade Zahl
kleiner, als die Zahl der beim Uebergang von $\alpha$ zu $\beta$ in der
Doppelreihe (7) gewonnenen Folge-Folgen.}

Von diesen beiden Sätzen, die nicht immer dieselbe obere Grenze für die
Zahl der Wurzeln geben, wird man den anwenden, der die niedrigste Grenze
giebt.
